\documentclass[10pt,UTF8]{article}


\usepackage{geometry}

\usepackage[scheme=plain]{ctex}

\geometry{a4paper,scale=0.9}

% \usepackage{times}  % 设置正文字体为 Adobe Times Roman
% \usepackage{mathptmx} % 只加载数学字体
% \renewcommand{\rmdefault}{cmr} % 恢复正文字体为 Computer Modern Roman

\usepackage{bm} % 数学粗体（避免斜体被取消）

% \usepackage{makecell}
% \usepackage{multirow} % 表格多行合并

\usepackage{booktabs} % 表格美化

% \usepackage{graphicx} %插入图片的宏包
% \usepackage{subfigure}
% \usepackage{float} %设置图片浮动位置的宏包

\usepackage{amsmath}
\usepackage{esint} % 提供更丰富的积分符号
\usepackage{amssymb}
\usepackage{mathtools}

\usepackage{physics}

% \usepackage{siunitx}
% % 关闭警告
% \ExplSyntaxOn
% \msg_redirect_name:nnn { siunitx } { physics-pkg } { none }
% \ExplSyntaxOff

% \usepackage{bussproofs} % 自然演绎推理树

% 交叉引用
% 关闭边框，设置颜色
\usepackage[colorlinks=true, linkcolor=darkgray, citecolor=teal,
urlcolor=blue]{hyperref}
% 使 \cref 和 \Cref 可用
\usepackage{cleveref}

\usepackage{bookmark}  % 生成 PDF 大纲

\allowdisplaybreaks[3]  % 允许换页


%%%%%%%%%%%%%%%%%%%%%%%% tcolorbox %%%%%%%%%%%%%%%%%%%%%%%%%%
\usepackage[most]{tcolorbox}

\newenvironment{tipBox}
{
  \begin{tcolorbox}[colback=green!5!white,colframe=green!75!black,parbox=false,before
  upper=\par,before lower=\par]}
  {
\end{tcolorbox}}

\newenvironment{conceptBox}
{
\begin{tcolorbox}[colback=blue!5!white,colframe=blue!75!black]}
  {
\end{tcolorbox}}

\newenvironment{theoremBoxNoIndent}
{
\begin{tcolorbox}[colback=yellow!5!white,colframe=yellow!75!black]}
  {
\end{tcolorbox}}

\newenvironment{definitionBox}
{
  \begin{tcolorbox}[colback=blue!5!white,colframe=blue!75!black,parbox=false,before
  upper=\par,before lower=\par]}
  {
\end{tcolorbox}}

\newenvironment{theoremBox}
{
  \begin{tcolorbox}[colback=yellow!5!white,colframe=yellow!75!black,parbox=false,before
  upper=\par,before lower=\par]}
  {
\end{tcolorbox}}
%%%%%%%%%%%%%%%%%%%%%%%%%%%%%%%%%%%%%%%%%%%%%%%%%%%%%%%%%%%%

%%%%%%%%%%%%%%%%%%%%% amsthm %%%%%%%%%%%%%%%%
\usepackage{amsthm}

\theoremstyle{definition}
\newtheorem{definition inner}{Definition}[section]
\newenvironment{definition}[1][]
{
  \begin{definitionBox}
  \begin{definition inner}[#1]\hfill\par}
    {
    \end{definition inner}
\end{definitionBox}}

\theoremstyle{definition}
\newtheorem{theorem inner}{Theorem}[section]
\newenvironment{theorem}[1][]
{
  \begin{theoremBox}
  \begin{theorem inner}[#1]\hfill\par}
    {
    \end{theorem inner}
\end{theoremBox}}

\theoremstyle{definition}
\newtheorem*{proof inner}{\textit{Proof}}
\renewenvironment{proof}[1][]
{
\begin{proof inner}[#1]\hfill\par}
  {\qed\vphantom{.}
\end{proof inner}}

\theoremstyle{plain}
\newtheorem{tip inner}{Tip}[section]
\newenvironment{tip}[1][]
{
  \begin{tipBox}
  \begin{tip inner}[#1]\hfill\par}
    {
    \end{tip inner}
\end{tipBox}}

\theoremstyle{remark}
\newtheorem*{remark}{\textit{Remark}}
%%%%%%%%%%%%%%%%%%%%%%%%%%%%%%%%%%%%%%%%%%%%%%


\newcommand{\iu}{\mathrm{i}}  % 虚数单位 i

\DeclareMathOperator{\Null}{Null}  % Null space


\begin{document}

\part*{Review of Matrices}

\section{}


\begin{definition}[Eigenvalues and Eigenvectors]
  For a square matrix \(A\), if there exists a non-zero vector \(\bm{v}\)
  and a scalar \(\lambda \) such that
  \[
    A\bm{v} = \lambda\bm{v},
  \]
  then \(\lambda \) is called an eigenvalue of \(A\) and
  \(\bm{v}\) is called an eigenvector corresponding to \(\lambda \).
\end{definition}

\begin{definition}[Multiplicity of Eigenvalues]
  Suppose \(\lambda \) is an eigenvalue of a matrix \(A\) and let the
  characteristic polynomial of \(A\) be \(p(x) = \det(A - x I)\).
  Then we have two types of multiplicity for \(\lambda \):
  \begin{itemize}
    \item \textbf{Algebraic multiplicity} of \(\lambda \):
      The number of times \(\lambda \) appears as a root of \(p(x)\).
    \item \textbf{Geometric multiplicity} of \(\lambda \):
      It is equivalent to any of the following (all of which are equal):
      \begin{itemize}
        \item The number of linearly independent eigenvectors
          corresponding to \(\lambda \).
        \item The dimension of the corresponding eigenspace:
          \( \dim(\{\bm{v} \mid A\bm{v} = \lambda\bm{v}\}) \).
        \item The nullity of \(A - \lambda I\):
          \( \dim(\Null(A - \lambda I)) \).
      \end{itemize}
  \end{itemize}
\end{definition}

\begin{theorem}
  For any eigenvalue \(\lambda \) of a square matrix \(A\),
  \begin{equation*}
    \text{geometric multiplicity of } \lambda
    \quad \leqslant \quad
    \text{algebraic multiplicity of } \lambda.
  \end{equation*}
\end{theorem}
\begin{proof}
  Suppose \(A \in \mathbb{R}^{n \times n}\) and \(\lambda \) is
  an eigenvalue of \(A\) with geometric multiplicity \(k\).
  Since \(\dim(\Null(A - \lambda I)) = k\),
  we can pick \(k\) orthonormal eigenvectors
  \( \{\bm{v}_1, \bm{v}_2, \ldots, \bm{v}_k\} \) from \(\Null(A - \lambda I)\),
  and then extend this set to an orthonormal basis of \(\mathbb{R}^n\)
  by adding \(n-k\) more vectors
  \( \{\bm{v}_{k+1}, \bm{v}_{k+2}, \ldots, \bm{v}_n\} \).
  Define the following matrices:
  \begin{equation*}
    V_1 =
    \begin{bmatrix}
      \vrule & & \vrule \\
      \bm{v}_1 & \cdots & \bm{v}_k \\
      \vrule & & \vrule
    \end{bmatrix}
    \qquad
    V_2 =
    \begin{bmatrix}
      \vrule & & \vrule \\
      \bm{v}_{k+1} & \cdots & \bm{v}_n \\
      \vrule & & \vrule
    \end{bmatrix}
    \qquad
    V =
    \begin{bmatrix}
      & & & & \\
      & \ldots V_1 & & \ldots V_2 & \\
      & & & &
    \end{bmatrix} =
    \begin{bmatrix}
      \vrule & & \vrule \\
      \bm{v}_{1} & \cdots & \bm{v}_n \\
      \vrule & & \vrule
    \end{bmatrix}
  \end{equation*}
  It is easy to know that
  \(V_1 \in \mathbb{R}^{n \times k} \) and
  \(V_2 \in \mathbb{R}^{n \times (n-k)}\),
  and \(V \in \mathbb{R}^{n \times n}\) is an orthogonal matrix.
  Then we have
  \begin{align*}
    A V_1 = \lambda V_1,
    \qquad\qquad
    V^\top A V &=
    \begin{bmatrix}
      & & \\
      & \ldots V_1^\top & \\
      & & \\
      & \ldots V_2^\top & \\
      & &
    \end{bmatrix} A
    \begin{bmatrix}
      & & & & \\
      & \ldots V_1 & & \ldots V_2 & \\
      & & & &
    \end{bmatrix} \\
    &=
    \begin{bmatrix}
      & & \\
      & \ldots (V_1^\top A V_1) & \ldots (V_1^\top A V_2) & \\
      & & \\
      & \ldots (V_2^\top A V_1) & \ldots (V_2^\top A V_2) & \\
      & &
    \end{bmatrix} \\
    &=
    \begin{bmatrix}
      & & \\
      & \ldots (\lambda I_k) & \ldots (V_1^\top A V_2) & \\
      & & \\
      & \ldots O_{n-k} & \ldots (V_2^\top A V_2) & \\
      & &
    \end{bmatrix}
    \in \mathbb{R}^{n \times n}
  \end{align*}
  Thus, the characteristic polynomial of \(A\) can be computed as
  \begin{align*}
    p(x) =
    \det(A - x I)
    &= \det(V^\top (A - x I) V ) \\
    &= \det(V^\top A V - x I) \\
    &= \det\left(
      \begin{bmatrix}
        & & \\
        & \ldots (\lambda I_k - x I_k) & \ldots (V_1^\top A V_2) & \\
        & & \\
        & \ldots O_{n-k} & \ldots (V_2^\top A V_2 - x I_{n-k}) & \\
        & &
      \end{bmatrix}
    \right) \\
    &= {(\lambda - x)}^k
    \cdot
    \det(V_2^\top A V_2 - x I_{n-k})
  \end{align*}
  Therefore, \(\lambda \) is a root of \(p(x)\) with multiplicity at
  least \(k\),
  which means the algebraic multiplicity of \(\lambda \) is at least \(k\).
\end{proof}


\end{document}
